\section{Orthogonal Complement of a Subspace}

\textbf{(a)}
% Let
% \[
%     \mathcal{V}^\perp=\{x\in\mathbb{R}^n\mid \langle x,y\rangle=0,\ \forall y\in V\}.
% \]
\begin{itemize}
    \item Zero in subset

    \(\langle 0,y\rangle=0^{\top}y=0\), so \(0\in \mathcal{V}^\perp\)
    \item Closure under addition
    
    Let \(x_1,x_2\in \mathcal{V}^\perp\), then for any \(y\in V\),
\[
\langle x_1, y \rangle = 0, \qquad 
\langle x_2, y \rangle = 0
\]

    Then
\[
    \langle x_1 + x_2,y\rangle
    =
    \langle x_1,y\rangle+\langle x_2,y\rangle
    =
    0
\]

Therefore \(x_1+x_2\in \mathcal{V}^\perp\)
    \item Closure under scalar multiplication
    
    Take any \(x\in \mathcal{V}^\perp\) and \(\alpha\in\mathbb{R}\)
\[
    \langle \alpha x,y\rangle
    =
    \alpha\langle x,y\rangle
    =
    \alpha x^{\top}y
    =
    0
\]
Therefore \(\alpha x \in \mathcal{V}^{\perp}\)

\textbf{Hence \(\mathcal{V}^\perp\) is a subspace of \(\mathbb{R}^n\)}

\end{itemize}

\textbf{(b)}
% Let
% \[
%     \mathcal{V}=
%     \begin{bmatrix}
%         v_1&v_2&\cdots&v_r
%     \end{bmatrix}
%     \in\mathbb{R}^{n\times r}.
% \]
% Then
\[
    \mathcal{V}=\operatorname{range}(V)
\]
% A vector \(x\) is in \(\mathcal{V}^\perp\) exactly when it is orthogonal to every
% column of \(\mathcal{V}\), i.e.,
% \[
%     \mathcal{V}^{\top} x=0.
% \]
% Therefore
\[
    \mathcal{V}^\perp=\operatorname{null}(\mathcal{V}^{\top})
\]

\textbf{(c)}

Let \(v\) be the orthogonal projection of \(x\) onto \(\mathcal{V}\), then
\[
    v^\perp=x-v
\]

By the definition of orthogonal projection, \(v\in \mathcal{V}\) and \(x-v\) is orthogonal to every vector in \(\mathcal{V}\).

Therefore,
\[
    v^\perp\in \mathcal{V}^\perp
\]

and
\[
    x=v+v^\perp
\]

To prove uniqueness, suppose
\[
    x=v_1+w_1=v_2+w_2
\]

where \(v_1,v_2\in \mathcal{V}\) and \(w_1,w_2\in \mathcal{V}^\perp\). Then
\(
    v_1-v_2=w_2-w_1
\)
\newline

The LHS belongs to \(\mathcal{V}\), while the RHS belongs to \(\mathcal{V}^\perp\). Hence this common vector belongs to
both $\mathcal{V}$ and $\mathcal{V}^\perp$: \(\mathcal{V}\cap \mathcal{V}^\perp\).
\newline

The only vector in this intersection is zero. If
\(z\in \mathcal{V}\cap \mathcal{V}^\perp\), then \(z\) is orthogonal to itself, so
\[
    z^\top z=\|z\|_2^2=0
\]
which implies \(z=0\). Therefore,
\(
    v_1=v_2
    \text{ and }
    w_1=w_2
\)

\textbf{So it is unique}.

\textbf{(d)}

Using \(V\) from part (b),
\[
    \dim \mathcal{V}=\operatorname{rank}(\mathcal{V}),
    \qquad
    \dim \mathcal{V}^\perp=\dim\operatorname{null}(\mathcal{V}^{\top}).
\]

By rank-nullity applied to \(\mathcal{V}^{\top}\),
\[
    \dim\operatorname{null}(\mathcal{V}^{\top})
    +
    \operatorname{rank}(\mathcal{V}^{\top})
    =
    n
\]

Since \(\operatorname{rank}(\mathcal{V}^{\top})=\operatorname{rank}(\mathcal{V})\),
we get
\[
    \dim \mathcal{V}^\perp+\dim \mathcal{V}=n
\]

\textbf{(e)}

Suppose \(\mathcal{V}\subseteq \mathcal{U}\), and let \(x\in \mathcal{U}^\perp\). Then \(x\) is
orthogonal to every vector in \(\mathcal{U}\). Since every vector in \(\mathcal{V}\) is also
in \(\mathcal{U}\), \(x\) is orthogonal to every vector in \(\mathcal{V}\). Thus
\[
    x\in \mathcal{V}^\perp
\]

Therefore
\[
    \mathcal{U}^\perp\subseteq \mathcal{V}^\perp
\]
